\documentclass[11pt,a4paper]{article}

% --- Essential Packages ---
\usepackage[utf8]{inputenc}
\usepackage{amsmath, amssymb, amsthm, mathtools} % For advanced math
\usepackage{graphicx}                            % For including figures
\usepackage{algorithm, algpseudocode}            % For pseudocode algorithms
\usepackage{listings}                            % For source code listings
\usepackage{xcolor}                              % For colors in code/listings
\usepackage{hyperref}                            % For clickable links

% --- Theorem and Proof Environments ---
\newtheorem{theorem}{Theorem}[section]
\newtheorem{lemma}[theorem]{Lemma}
\theoremstyle{definition}
\newtheorem{definition}[theorem]{Definition}

% --- Code Listing Styling ---
\lstset{
    basicstyle=\ttfamily\small,
    breaklines=true,
    frame=single,
    numberstyle=\tiny\color{gray},
    numbers=left,
    keywordstyle=\color{blue},
    commentstyle=\color{green!60!black},
    stringstyle=\color{red}
}

\title{Computer Science Research Document: Math, Algorithms, and Figures}
\author{Your Name}
\date{\today}

\begin{document}

\maketitle

\begin{abstract}
This document demonstrates how to typeset mathematics, pseudocode, program source code, and figures in LaTeX.
\end{center}
\end{abstract}

\section{Introduction and Math Notation}
LaTeX provides robust tools for mathematical expressions. You can write inline equations like $f(n) = O(n \log n)$ or use display math:

\begin{equation}
    T(n) = 2T\left(\frac{n}{2}\right) + \Theta(n)
    \label{eq:master_theorem}
\end{equation}

\begin{equation}
    \int_{0}^{\infty } x\,\mathrm{d}x
    \label{eq:master_theorem}
\end{equation}

By the Master Theorem, the solution to Equation~(\ref{eq:master_theorem}) is $T(n) = \Theta(n \log n)$.

\section{Algorithms}
You can write structured pseudocode using the [algpseudocode](https://www.overleaf.com/learn/latex/Algorithms) package:

\begin{algorithm}
\caption{Binary Search Algorithm}
\begin{algorithmic}[1]
\Require Sorted array $A$ of size $n$, target value $t$
\Ensure Index of $t$ if found, else $-1$
\State $low \gets 0$
\State $high \gets n - 1$
\While{$low \le high$}
    \State $mid \gets \lfloor (low + high) / 2 \rfloor$
    \If{$A[mid] = t$}
        \State \textbf{return} $mid$
    \ElsIf{$A[mid] < t$}
        \State $low \gets mid + 1$
    \Else
        \State $high \gets mid - 1$
    \EndIf
\EndWhile
\State \textbf{return} $-1$
\end{algorithmic}
\label{alg:binary_search}
\end{algorithm}

\section{Computer Science Source Code}
To include actual programming code, use the `listings` package:

\begin{lstlisting}[language=Python, caption={Python implementation of Binary Search}]
def binary_search(A, t):
    low, high = 0, len(A) - 1
    while low <= high:
        mid = (low + high) // 2
        if A[mid] == t:
            return mid
        elif A[mid] < t:
            low = mid + 1
        else:
            high = mid - 1
    return -1
\end{lstlisting}

\section{Figures and Graphics}
Figures are handled via the `graphicx` package. Place your image file (e.g., `diagram.png`) in your project folder:

\begin{figure}[htbp]
    \centering
    \includegraphics[width=0.6\textwidth]{diagram.png}
    \caption{An architectural overview of the system flow.}
    \label{fig:architecture}
\end{figure}

As shown in Figure~\ref{fig:architecture}, the data flows directly from the input layer to the processing node.

\end{document}

